   \subsection{Stakeholders por parte del cliente y por parte de la empresa}
De acuerdo con el proceso de Recopilar Requisitos, el éxito del proyecto depende de la participación activa de estos interesados para transformar sus necesidades en requisitos medibles.

\begin{table}[H]
\footnotesize
\centering
\begin{tabulary}{\textwidth}{@{}LLLLL@{}}
\toprule
\textbf{Parte} & \textbf{Nombre/CC} & \textbf{Cargo} & \textbf{Celular} & \textbf{Email} \\ \midrule
\textbf{Cliente (Estado)} & Luis Alejandro Zambrano Ruiz / 1111111111 & Director TI (Mintransporte) & (+57 601) 3240800 & alezam30@hotmail.com \\
\textbf{Cliente (Estado)} & David Enrique Mayorga Peláez / 1111111111 & Coord. Nacional RUNT & (57+1) 3240800 Ext. 1045. & servicioalciudadano@mintransporte.gov.co \\
\textbf{Empresa (Software)} & Jorge Luis Araque / 1111111111 & Director de Proyecto (PM) & (+57) 311 3315381 & jorge.araque@utp.edu.co \\
\textbf{Empresa (Software)} & Juan Esteban Cardona / 1113858851 & Arquitecto de Software & (+57) 301 2213085 & juan.cardona@utp.edu.co \\
\textbf{Empresa (Software)} & Jhoan Esteban Corrales / 1113859081 & Arquitecto de Software & (+57) 311 3788820 & jhoan.corrales@utp.edu.co \\
\textbf{Empresa (Software)} & Juan Camilo Velasquez / 1111111111 & Arquitecto de Software & (+57) 313 6992715 & juancamilo.velasquez@utp.edu.co \\
\textbf{Gremio Interesado} & Gremio Carga / N/A & Representante Colfecar & (+57) 320 303 8830. & info@colfecar.org.co \\ \bottomrule
\end{tabulary}
\end{table}